\documentclass[instructions.tex]{subfiles}
\begin{document}
 
\chapter{De la médisance, de l’envie et des jugemens téméraires.}
 
\primo Nous devons tâcher de conserver la réputation d’autrui, comme nous voudrions qu’on conservât la nôtre, et plus craindre d’attenter à son honneur qu’à ses biens. Vous ne devriez pas souiller vos mains du bien d’autrui; mais vous êtes bien plus coupable, lorsque vous souillez votre langue par la médisance. Le médisant est ( toute proportion gardée ) plus criminel qu’un voleur, parce qu’en ôtant au prochain sa réputation, il lui ôte quelque chose de plus précieux que ses biens\footnote{Melius est nomen bonum, quàm divitiæ multæ.}.
 
Estimons tout le monde. Tel nous paraît vicieux, qui, devant Dieu, vaut mieux que nous: il sera peut-être un jour dans le ciel par sa pénitence, tandis que nous en serons exclus nous-mêmes, si nous manquons de charité. Prendre plaisir à remarquer les vices d’autrui, s’entretenir de ses défauts, et jamais de ses vertus; critiquer ses paroles, ses actions, ses manières, ce n’est pas aimer le prochain, c’est le traiter en ennemi. Quelle bizarrerie dans l’homme ! A peine en trouve-t-on un qui borne ses vues sur soi-même. On est aveugle sur ses propres défauts, et éclairé sur les défauts d’autrui. Vous voyez, dit l’Évangile, un fétu dans l’œil de votre frère, et vous ne voyez pas la poutre qui est dans le vôtre. Mais pourquoi tant de réflexions sur les autres et si peu sur soi ? N’y a-t-il pas dans nous-mêmes de quoi nous occuper et nous humilier, sans porter nos vues sur les autres ? Pensons à nos défauts dont nous répondrons à Dieu, et laissons les défauts des autres, dont nous ne sommes pas responsables.
 
Comprenons combien nous blessons la charité en parlant mal d’autrui. Une parole qu’on a dite de nous ou de notre famille, nous perce le cœur et nous afflige; pourquoi prenons-nous plaisir d’affliger et de flétrir les autres ? Cette habitude de parler mal du prochain ne peut venir que d’un grand fonds d’imperfection et d’orgueil. On remarque en effet que ce sont les plus imparfaits, ceux sur lesquels il y a le plus à reprendre, qui parlent beaucoup d’autrui.
 
On parle mal des autres, on les condamne, parce qu’on est soi-même méchant et souvent plus méchant. Vous vous condamnez, dit saint Paul, lorsque vous condamnez les autres, car vous êtes disposé vous-même à faire le mal que vous jugez d’autrui; et peut-être faites-vous de plus grandes fautes. La bouche parle de l’abondance du cœur, dit Jésus-Christ; il y a donc dans votre cœur des désordres en abondance, puisqu’il se répand si souvent contre le prochain.
 
On éprouve tous les jours ce que dit saint Jacques, que la langue, une des plus petites parties du corps, cause les plus grands maux. Elle porte le trouble dans la société, dans les familles, parmi les amis et dans les mariages. Apprenons donc à parler, apprenons à écouter et à nous taire; voilà la marque d’un homme parfait\footnote{Si quis in verbo non offendit, hic perfectus est vir. Jac. 3.}.
 
Il est néanmoins quelquefois permis de faire connoître avec prudence les fautes et les défauts d’autrui, dans le cas où le bien public, l’utilité de celui de qui on parle, ou l’avantage de celui à qui l’on parle, le demandent. Hors ces cas, vous devez réparer la réputation et les suites, si la chose est de conséquence, et si elle est ignorée dans le lieu où vous l’avez dite,
 
S’il est rare de trouver de la discrétion et de la charité dans ceux qui parlent, il n’est pas moins rare d’en trouver dans ceux qui écoutent. Ceux-là ont le démon sur la langue, dit saint Bernard, et ceux-ci dans l’oreille. Faites taire le médisant, si vous avez l’autorité; du moins ne consentez pas, n’ajoutez pas même foi à tout le mal qu’on dit d’autrui. C’est une légèreté de croire d’abord tout ce que l’on entend, dit l’Écriture. Il arrive rarement que deux personnes racontent une chose de la même façon; on a donc sujet de douter de ce qu’on ne sait que par ouï dire. Souvent le premier qui l’a dite est dans l’erreur, et accuse faussement un innocent; et quand il aurait dit la vérité, il est toujours plus avantageux de ne pas croire ou de suspendre son jugement. Nous ne devons penser mal que de nous-mêmes.
 
\secundo Nous n’aimons pas les autres, nous en jugeons et nous en parlons mal; pourquoi ? parce que nous n’aimons que nous-mêmes, et parce que nous sommes envieux. Ce ne sont pas toujours leurs défauts que nous haïssons, c’est leur mérite et leur fortune. La jalousie nous aveugle et nous prévient. Aussitôt qu’un voisin, un parent, un concurrent, a plus d’esprit, plus de talent, plus de bien, on ne peut le souffrir. On donne de malignes interprétations à tout ce qu’il dit, à tout ce qu’il fait; tout déplaît en lui jusqu’à ses belles qualités. Quelle bassesse de cœur ! Pourquoi les pharisiens ne pouvaient-ils souffrir le Sauveur ? C’est parce qu’ils étaient possédés par l’envie. Le démon ne peut souffrir les dons de Dieu dans les autres; il hait les hommes et voudrait les perdre, parce que Dieu les bénit et les aime. Telle est la passion d’un envieux, Passion maligne: il ne veut que le mal, il n’aime les autres que quand ils sont malheureux. Passion aveugle, qu’on déteste dans les autres, et qu’on ne veut pas reconnoître dans soi-même. Passion cruelle, qui porte avec elle son supplice. Passion bizarre, qui s’afflige de ce qui réjouit les autres, et qui fait son malheur du bonheur d’autrui. O envieux ! pourquoi vous rongez-vous le cœur sur la prospérité de vos frères ? Soyez contens de ce que vous avez, et de ce que Dieu veut que vous ayez. Laissez les autres en paix; ne nourrissez pas dans votre esprit des monstres qui vous déchirent l’âme et qui vous damnent.
 
\tertio Si l’on se prévient contre le prochain, si l’on en parle mal, c’est parce qu’on en juge mal; mais reconnoissons ici l’injustice de nos jugemens. Nous ne nous connoissons pas nous-mêmes, et nous prétendons connoître les autres. Lorsque nous regardons notre conscience, nous ne pouvons souvent en démêler les motifs, ni juger si nos intentions sont bonnes ou mauvaises, ou si nous avons consenti à une pensée ou non. Comment donc osons-nous juger des intentions et des pensées d’une personne dont nous ne connoissons pas la conscience et le cœur ?
 
C’est une grande témérité de s’ériger en juge des autres et en censeur des sentimens de leur âme, qui ne sont connus que de Dieu. Juger de la sorte, dit saint Jean Climaque, c’est porter l’insolence jusque sur le trône de la Divinité, et attenter aux droits de Jésus-Christ le juge souverain\footnote{Judicare est impudens direptio divinæ auctoritatis. Grad. 10.}. Nous serons jugés comme nous aurons jugé les autres; et si Dieu nous juge aussi impitoyablement que nous jugeons nos frères, notre condamnation est assurée.
 
Le propre caractère de la charité est d’être bon, de penser du bien de tous. Quand on aime véritablement, on trouve toujours, même dans ceux qui ont des défauts, de quoi les louer, ou du moins les excuser. Il n’y a que les démons et les réprouvés qui soient indignes de notre charité et de nos bienfaits.
 
\section{Résolutions}
 
\primo Je regarderai toujours la réputation du prochain comme sa propriété, à laquelle je ne dois toucher ni par médisance, ni par raillerie, ni par jugement téméraire. 

\secundo Je ne lui porterai point envie; au contraire, je me réjouirai de son bonheur et de son mérite, et je louerai volontiers en lui, en son absence, ce que j’y connoîtrai de bon. 

\tertio Au lieu de prendre part à la médisance, quand j’en entendrai faire, j’imposerai silence au détracteur, si j’en ai l’autorité: je détournerai adroitement le discours sur un autre objet, je montrerai un visage triste. 

\prayer{Mon Dieu, faites-moi la grâce de pouvoir dire désormais, comme un de vos prophètes: J’étais le fléau de celui qui médisait\footnote{iDetrahentem secretò proximo suo, hunc persequebar. Ps.100.5.}.}
 
\end{document}
